% SHARD-ID: AQM-92-HAAH-DATA-SCHEME-MODULES
% SHARD-TITLE: Haah Data over Schemes: Permissive Module Grammar
% SHARD-SUMMARY: Defines a broad pre-Haah datum over a commutative ring and its scheme interpretation over \(\Spec A\).
% SHARD-SUMMARY: Separates algebraic module constraints, adjoint/excitation maps, and defect quotients from Weyl, Hilbert, Hamiltonian, locality, and code layers.
% SHARD-SUMMARY: Tests the grammar on fields, product rings, dual numbers, \(\mathbb Z/6\mathbb Z\), and \(\mathbb F_3[t]/(t^n-1)\).
% SHARD-KEYWORDS: Haah, scheme module, Spec, Laurent module, symplectic form, excitation map, defect module, nilpotent

\section{Haah Data over Schemes: Permissive Module Grammar}
\label{sec:haah-data-scheme-modules}

\subsection{Status and Source Anchors}

\textbf{Status: local definitions and proposals, with source-local Haah
inputs.}  This shard implements convention \texttt{(ba)}: a pre-Haah datum is
only an algebraic control package, a scheme Haah datum is its
\(\Spec(A)\)-interpretation through sheaves, and neither layer by itself gives
Pauli operators, Hilbert spaces, locality, Hamiltonians, or codes
(\path{CONVENTIONS.md}, lines 2462--2505).  The affine-scheme and module-sheaf
background is Stacks: \(\Spec(R)\), \(D(f)\), and residue fields are in
\path{references/algebraic_geometry/stacks_project_algebra.tex}, lines
273--278 and 2972--3023; associated module sheaves, global sections,
localizations, stalks, and exactness of \(M\mapsto\widetilde M\) are in
\path{references/algebraic_geometry/stacks_project_schemes.tex}, lines
664--718.

The Haah source-local inputs are fixed by the lattice-stabilizer manifest
\path{references/lattice_stabilizer_codes/SOURCES.md}, lines 17--43.  In the
free-module source, Haah encodes translation-invariant commuting Pauli
Hamiltonians by maps between free modules over a Laurent/translation group
algebra
(\path{references/lattice_stabilizer_codes/haah_1204_1063_source/alg-code-hamiltonian-v6.tex},
lines 168--183).  The Pauli module over \(R=\mathbb F_2[\Lambda]\), the
antipode, trace, commutation form, \(P=R^{2q}\), and
\(\sigma^\dagger\lambda_q\sigma=0\) criterion are source-local in the same
file, lines 411--525 and 528--544.  The generator-label map
\(\sigma:G\to P\) and excitation map \(\epsilon=\sigma^\dagger\lambda_q\) are
also in Haah's thesis
(\path{references/lattice_stabilizer_codes/haah_1305_6973_source/alg-theory.tex},
lines 320--337).  The complex condition
\(\epsilon\circ\sigma=0\) and \(\operatorname{im}\sigma\subseteq\ker\epsilon\)
are source-local in
\path{references/lattice_stabilizer_codes/haah_1204_1063_source/alg-code-hamiltonian-v6.tex},
lines 656--667.  AQM-91 records the finite periodic bridge and warns that
Haah's translation Laurent ring is not automatically the coordinate ring of an
algebraic torus (AQM-91, lines 41--91 and 165--200).

\subsection{The Source-Local Grammar to Generalize}

Haah's source-local data have five algebraic pieces before any scheme
generalization is attempted.  First, there is a coefficient ring, in the source
case a translation group algebra.  Second, the Pauli label module \(P\) is free
over that ring.  Third, a generator-label module \(G\) maps to \(P\) by
\(\sigma:G\to P\).  Fourth, a Laurent symplectic form \(\lambda_q\), together
with the dagger induced by transpose and the antipode, gives the excitation
map \(\epsilon=\sigma^\dagger\lambda_q\).  Fifth, isotropy or commutation is
the equation \(\sigma^\dagger\lambda_q\sigma=0\).  The preceding paragraph
cites the local source lines for each item.

\textbf{Local proposal.}  The general lesson kept here is not ``every
\(\Spec(A)\) is a Pauli lattice''.  The lesson is the module grammar
\[
  G \xrightarrow{\sigma} P,\qquad
  P \xrightarrow{\epsilon} G^\vee,\qquad
  \epsilon\sigma=0,
\]
plus the habit of recording the quotient \(\ker\epsilon/\operatorname{im}\sigma\)
when it is defined.  All translation, finite support, Pauli, Hilbert, and
Hamiltonian meanings are deliberately stripped away unless a later layer adds
the extra data named in convention \texttt{(ba)}.

\subsection{Pre-Haah Data over a Ring}

\paragraph{Definition D1: permissive pre-Haah datum.}
Let \(A\) be a commutative ring.  A \emph{pre-Haah datum over \(A\)} is a tuple
\[
  \mathcal D=(G,P,\sigma,\lambda_{\rm opt})
\]
where \(G\) and \(P\) are specified \(A\)-modules, \(\sigma:G\to P\) is an
\(A\)-linear map, and \(\lambda_{\rm opt}\) is either absent or is a specified
\(A\)-linear map
\[
  \lambda:P\to P^\vee,\qquad P^\vee=\operatorname{Hom}_A(P,A).
\]
When \(\lambda\) is absent, the datum records only the generated submodule
\(\operatorname{im}\sigma\subseteq P\).  It has no excitation map and no defect
quotient.

\paragraph{Definition D2: formed and isotropic datum.}
If \(\lambda\) is present, define the formal excitation map
\[
  \epsilon=\sigma^\vee\circ\lambda:P\longrightarrow G^\vee.
\]
This is a local definition in the category of \(A\)-modules; it is the
identity-involution analogue of Haah's source-local
\(\sigma^\dagger\lambda_q\).  The formed datum is \emph{isotropic} if
\[
  \epsilon\sigma=\sigma^\vee\lambda\sigma=0.
\]
Only in the formed isotropic case do we define the formal logical-defect module
\[
  \mathcal L_{\rm alg}(\mathcal D)
  =
  \ker(\epsilon)/\operatorname{im}(\sigma).
\]
This quotient is an algebraic module diagnostic.  It is not a physical logical
operator group unless a finite Weyl-Haah or translation Haah datum supplies the
missing representation and locality data, as required by
\path{CONVENTIONS.md}, lines 2493--2505.

\paragraph{Definition D3: perfection and finite local bases.}
The word \emph{symplectic} is reserved here for a chosen alternating
\(\lambda:P\to P^\vee\) that is an isomorphism.  Matrix adjoints, local ranks,
residue-fibre nondegeneracy, and base-changed duals require finite locally free
or finite projective hypotheses.  Stacks defines finite locally free modules in
\path{references/algebraic_geometry/stacks_project_algebra.tex}, lines
19004--19017, and records the equivalence with finite projective conditions in
lines 19027--19053.  The dual of a finite locally free sheaf is finite locally
free, and tensoring with one is exact
(\path{references/algebraic_geometry/stacks_project_cohomology.tex}, lines
13836--13845).  The Hom-tensor comparison used when duals are moved through
base change is guaranteed for finite projective modules in
\path{references/algebraic_geometry/stacks_project_algebra.tex}, lines
19351--19357.  Without these hypotheses, \(\epsilon\) may still be a formal
module map, but fibrewise matrix claims are not made here.

\subsection{Scheme Haah Data over \texorpdfstring{\(\Spec A\)}{Spec A}}

\paragraph{Definition D4: scheme interpretation.}
Let \(X=\Spec(A)\).  A \emph{scheme Haah datum} is a pre-Haah datum over \(A\)
viewed through the quasi-coherent sheaves
\[
  \widetilde G,\qquad \widetilde P,\qquad
  \widetilde\sigma:\widetilde G\to\widetilde P,
\]
and, when present, \(\widetilde\lambda:\widetilde P\to\widetilde{P^\vee}\).
On a standard open \(D(f)\), the restricted modules are \(G_f\) and \(P_f\);
at \(x=\mathfrak p\), the stalks are \(G_{\mathfrak p}\) and
\(P_{\mathfrak p}\) by the Stacks affine-module-sheaf result cited above.
The residue fibre used in this shard is the local definition
\[
  P(x)=P_{\mathfrak p}\otimes_{A_{\mathfrak p}}\kappa(\mathfrak p),
  \qquad
  G(x)=G_{\mathfrak p}\otimes_{A_{\mathfrak p}}\kappa(\mathfrak p).
\]

The permissive scheme layer allows quasi-coherent modules when the intended
operation is restriction, localization, stalks, support language, or residue
fibres.  It asks for coherent modules when a finite-equation or finite-support
scheme-theoretic statement is intended: on locally Noetherian schemes,
coherent is equivalent to quasi-coherent finite type/finitely presented in the
sense recorded in
\path{references/algebraic_geometry/stacks_project_coherent.tex}, lines
2224--2242, and quasi-coherent submodules and quotients of coherent modules
remain coherent by lines 2328--2344.  It asks for finite locally free modules
when the operation needs local bases, ranks, perfect forms, duals compatible
with fibres, or a matrix \(\sigma^\vee\lambda\sigma\) whose entries have stable
local meaning.

\subsection{Restrictive Ladder}

\begin{center}
\small\setlength{\tabcolsep}{3pt}
\begin{tabular}{p{0.22\linewidth}p{0.33\linewidth}p{0.34\linewidth}}
\textbf{Layer} & \textbf{Added data} & \textbf{Allowed claim} \\
\hline
Pre-Haah & \(A\)-modules \(G,P\), map \(\sigma\) &
\(\operatorname{im}\sigma\subseteq P\) is an algebraic submodule. \\
Formed pre-Haah & \(\lambda:P\to P^\vee\) &
\(\epsilon=\sigma^\vee\lambda\), isotropy, and
\(\ker\epsilon/\operatorname{im}\sigma\) if isotropic. \\
Scheme Haah & sheaves on \(X=\Spec A\) &
restriction, localization, stalks, residue fibres, support, and base-change
questions. \\
Finite locally free & finite projective/local-basis hypotheses &
matrix adjoints, perfect forms, fibrewise ranks, and dual-compatible
comparisons. \\
Finite Weyl or translation & generating characters, finite self-duality, group
algebra, antipode, trace, finite support &
Weyl labels or Haah translation-stabilizer interpretations; these are outside
this shard. \\
\end{tabular}
\end{center}

\subsection{Examples at the Scheme-Module Layer}

\paragraph{Field \(k\).}
\textbf{Local derivation.}  A field has exactly one prime ideal, \((0)\), so
\(\Spec(k)\) has one point.  For \(k=\mathbb F_p\), AQM-22 gives this
calculation and separates \(\Spec(\mathbb F_p)\) from the \(p\)-element
underlying set (AQM-22, lines 48--69).  A finite locally free pre-Haah datum is
just finite-dimensional \(k\)-vector spaces \(G\) and \(P\), a matrix
\(\sigma\), and optionally an alternating matrix \(\lambda\).  Then
\(\epsilon=\sigma^T\lambda\) and
\(\mathcal L_{\rm alg}=\ker\epsilon/\operatorname{im}\sigma\) are vector-space
objects if \(\sigma^T\lambda\sigma=0\).  The finite-field Weyl qudit statement
belongs to AQM-22/AQM-81, not to this bare module example.

\paragraph{Product field ring \(k^n\).}
AQM-40 proves that \(R=k^n\) has points
\(\mathfrak m_i=\ker(\operatorname{pr}_i)\), residues \(k\), and singleton
opens \(D(e_i)\) (AQM-40, lines 51--97); convention \texttt{(ad)} records that
\(\Spec(k^n)\) has \(n\), not \(q^n\), points for \(k=\mathbb F_q\)
(\path{CONVENTIONS.md}, lines 1012--1035).  \textbf{Local derivation:} since
the idempotents \(e_i\) are orthogonal and sum to \(1\), every \(R\)-module
\(M\) decomposes as \(\bigoplus_i e_iM\).  Thus a scheme Haah datum over
\(k^n\) is componentwise \(n\) field-level data after applying the idempotents.
AQM-40, lines 232--277, warns that stabilizer codes are later choices of
isotropic labels and phases; the product ring itself supplies no incidence
complex or lattice locality.

\paragraph{Dual numbers \(k[\varepsilon]/(\varepsilon^2)\).}
Stacks defines dual numbers as the \(k\)-module with basis \(1,\varepsilon\)
and \(\varepsilon^2=0\)
(\path{references/algebraic_geometry/stacks_project_varieties.tex}, lines
2918--2924).  \textbf{Local derivation:} every prime contains the nilpotent
\(\varepsilon\), so the topological spectrum is the one point of the quotient
\(k\), while the structure sheaf still remembers the nilpotent direction.
For a free module datum, write
\[
  \sigma=\sigma_0+\varepsilon\sigma_1,\qquad
  \lambda=\lambda_0+\varepsilon\lambda_1 .
\]
The isotropy equation expands to the zeroth-order condition
\(\sigma_0^T\lambda_0\sigma_0=0\) and the first-order condition
\[
  \sigma_1^T\lambda_0\sigma_0+
  \sigma_0^T\lambda_1\sigma_0+
  \sigma_0^T\lambda_0\sigma_1=0 .
\]
This is an algebraic deformation calculation only.  For
\(\mathbb F_3[\varepsilon]/(\varepsilon^2)\), AQM-82 records that the reduced
residue-site workflow has one \(\mathbb F_3\) site, while the thickened
Artin-Weyl layer needs the separate AQM-36 top-coefficient generating
character (AQM-82, lines 128--156).

\paragraph{\(\mathbb Z/6\mathbb Z\).}
AQM-23 proves the Chinese-remainder isomorphism
\(\mathbb Z/6\mathbb Z\cong\mathbb F_2\times\mathbb F_3\), computes the two
prime ideals \((2)\) and \((3)\), and gives residues \(\mathbb F_2\) and
\(\mathbb F_3\) (AQM-23, lines 57--100).  Therefore a module-level Haah datum
over \(\mathbb Z/6\mathbb Z\) splits into its two residue/product components
after the Chinese-remainder idempotents are chosen.  The formed equations
\(\epsilon\sigma=0\) and the algebraic defect module are componentwise
equations over the two fields.  AQM-23, lines 102--225, supplies a separate
finite Weyl residue construction for this ring; this shard uses none of that
as a consequence of the bare pre-Haah datum.

\paragraph{\(\mathbb F_3[t]/(t^n-1)\).}
AQM-35 writes \(n=3^s m\), \(3\nmid m\), proves
\[
  t^n-1=(t^m-1)^{3^s},
\]
and records that the closed residue Weyl sites are the irreducible factors of
the squarefree part \(t^m-1\) (AQM-35, lines 27--54).  At the present layer,
take \(A=R_n=\mathbb F_3[t]/(t^n-1)\) and choose modules, \(\sigma\), and
possibly \(\lambda\) over the full nonreduced ring.  The residue fibres of the
scheme Haah datum see the quotients \(K_\pi=\mathbb F_3[t]/(\pi)\), while
\(\sigma\), \(\lambda\), \(\ker\epsilon\), and
\(\mathcal L_{\rm alg}\) may still contain nilpotent information over \(R_n\).
AQM-35, lines 241--246, states that nilpotents are invisible to the reduced
residue-qudit layer, and AQM-37, lines 34--65 and 171--174, treats the
separate thickened Weyl layer.  No periodic Haah lattice is obtained here
unless a translation quotient, antipode, trace, and finite-support bridge like
AQM-91 is added.

\subsection{What This Layer Does Not Yet Give}

\begin{itemize}
\item It does not turn a bare \(\Spec(A)\) into sites of a lattice, a Pauli module, a Weyl representation, a Hilbert space, a Hamiltonian, locality, or a quantum code.
\item It does not supply Haah's dagger unless the needed duality, finite locally free hypotheses, and, for translation data, antipode/trace conventions are named.
\item It does not assert exactness \(\ker\epsilon=\operatorname{im}\sigma\); outside source-local Haah examples or checked finite computations, \(\ker\epsilon/\operatorname{im}\sigma\) is only a module to study.
\item It does not quantize arbitrary finite rings.  Finite Weyl-Haah data need finite self-duality, residue-field labels, or a certified Frobenius-ring generating character, exactly as convention \texttt{(ba)} requires.
\item It does not identify coordinate rings with translation group algebras.  That bridge is extra data, and AQM-91 is the local model for making it finite and explicit.
\end{itemize}
